% GENERATED FILE. Edit canonical lessons, then run npm run generate:books.
% canonical-source-hash: fnv1a64:88dc0884546cb5f5
% canonical-lessons: SA-C196-tamram, SA-C196-pasanah, SA-C196-ratnakarah, SA-C196-sarit, SA-C196-guha, SA-R196-first-pass-words-from-chapters-191-193, SA-R196-second-pass-words-from-chapters-194-196

\chapter{Copper, Rock, Ocean, Stream, Cave}
\label{ch:sa-copper-rock-ocean-stream-cave}
\hlchaptermodality{196}

\begin{chapteropening}
\textbf{By the end of this chapter:} \emph{I can say copper, a rock, the ocean, a stream and a cave.}

Complete the last lesson of chapter 196: I can say copper, a rock, the ocean, a stream and a cave.
\end{chapteropening}

\section[\texorpdfstring{\sk{ताम्रम्} (tāmram)}{ताम्रम् (tāmram)}]{\sk{ताम्रम्} (tāmram) — copper}

\label{lesson:SA-C196-tamram}



\begin{quote}
Before the new one: say the Sanskrit for a province, then the Sanskrit for a boundary.
\end{quote}

\subsection*{You'll want to know: \sk{ताम्रम्}}
\textbf{\sk{ताम्रम्}} — \emph{tāmram} — \textquotedblleft{}copper\textquotedblright{}.

\begin{grammarlens}[title={What you've built}]
One more word for this chapter.
\end{grammarlens}

\subsection*{Your turn}
\begin{itemize}
\raggedright
  \item \emph{Say it:} \emph{tāmram}
  \item \emph{Say it:} \emph{tāmram}, once more
  \item \emph{Say it:} say \emph{tāmram}
  \item \emph{Recall:} say the Sanskrit for a province, then the Sanskrit for a boundary, then say \emph{tāmram} again
\end{itemize}

\begin{tcolorbox}[breakable,colback=teal!4,colframe=teal!35!black,title={Before you move on}]
What does \sk{ताम्रम्} mean? (\textquotedblleft{}Copper\textquotedblright{}.) Say it once more.
\end{tcolorbox}
\section[\texorpdfstring{\sk{पाषाणः} (pāṣāṇaḥ)}{पाषाणः (pāṣāṇaḥ)}]{\sk{पाषाणः} (pāṣāṇaḥ) — a rock, a stone}

\label{lesson:SA-C196-pasanah}



\begin{quote}
Before the new one: say the Sanskrit for a boundary, then the Sanskrit for copper.
\end{quote}

\subsection*{You'll want to know: \sk{पाषाणः}}
\textbf{\sk{पाषाणः}} — \emph{pāṣāṇaḥ} — \textquotedblleft{}a rock, a stone\textquotedblright{}.

\begin{grammarlens}[title={What you've built}]
One more word for this chapter.
\end{grammarlens}

\subsection*{Your turn}
\begin{itemize}
\raggedright
  \item \emph{Say it:} \emph{pāṣāṇaḥ}
  \item \emph{Say it:} \emph{pāṣāṇaḥ}, once more
  \item \emph{Say it:} say \emph{pāṣāṇaḥ}
  \item \emph{Recall:} say the Sanskrit for a boundary, then the Sanskrit for copper, then say \emph{pāṣāṇaḥ} again
\end{itemize}

\begin{tcolorbox}[breakable,colback=teal!4,colframe=teal!35!black,title={Before you move on}]
What does \sk{पाषाणः} mean? (\textquotedblleft{}A rock, a stone\textquotedblright{}.) Say it once more.
\end{tcolorbox}
\section[\texorpdfstring{\sk{रत्नाकरः} (ratnākaraḥ)}{रत्नाकरः (ratnākaraḥ)}]{\sk{रत्नाकरः} (ratnākaraḥ) — the ocean}

\label{lesson:SA-C196-ratnakarah}



\begin{quote}
Before the new one: say the Sanskrit for copper, then the Sanskrit for a rock.
\end{quote}

\subsection*{You'll want to know: \sk{रत्नाकरः}}
\textbf{\sk{रत्नाकरः}} — \emph{ratnākaraḥ} — \textquotedblleft{}the ocean\textquotedblright{}.

\begin{grammarlens}[title={What you've built}]
One more word for this chapter.
\end{grammarlens}

\subsection*{Your turn}
\begin{itemize}
\raggedright
  \item \emph{Say it:} \emph{ratnākaraḥ}
  \item \emph{Say it:} \emph{ratnākaraḥ}, once more
  \item \emph{Say it:} say \emph{ratnākaraḥ}
  \item \emph{Recall:} say the Sanskrit for copper, then the Sanskrit for a rock, then say \emph{ratnākaraḥ} again
\end{itemize}

\begin{tcolorbox}[breakable,colback=teal!4,colframe=teal!35!black,title={Before you move on}]
What does \sk{रत्नाकरः} mean? (\textquotedblleft{}The ocean\textquotedblright{}.) Say it once more.
\end{tcolorbox}
\section[\texorpdfstring{\sk{सरित्} (sarit)}{सरित् (sarit)}]{\sk{सरित्} (sarit) — a stream, a river}

\label{lesson:SA-C196-sarit}



\begin{quote}
Before the new one: say the Sanskrit for a rock, then the Sanskrit for the ocean.
\end{quote}

\subsection*{You'll want to know: \sk{सरित्}}
\textbf{\sk{सरित्}} — \emph{sarit} — \textquotedblleft{}a stream, a river\textquotedblright{}.

\begin{grammarlens}[title={What you've built}]
One more word for this chapter.
\end{grammarlens}

\subsection*{Your turn}
\begin{itemize}
\raggedright
  \item \emph{Say it:} \emph{sarit}
  \item \emph{Say it:} \emph{sarit}, once more
  \item \emph{Say it:} say \emph{sarit}
  \item \emph{Recall:} say the Sanskrit for a rock, then the Sanskrit for the ocean, then say \emph{sarit} again
\end{itemize}

\begin{tcolorbox}[breakable,colback=teal!4,colframe=teal!35!black,title={Before you move on}]
What does \sk{सरित्} mean? (\textquotedblleft{}A stream, a river\textquotedblright{}.) Say it once more.
\end{tcolorbox}
\section[\texorpdfstring{\sk{गुहा} (guhā)}{गुहा (guhā)}]{\sk{गुहा} (guhā) — a cave}

\label{lesson:SA-C196-guha}



\begin{quote}
Before the new one: say the Sanskrit for the ocean, then the Sanskrit for a stream.
\end{quote}

\subsection*{You'll want to know: \sk{गुहा}}
\textbf{\sk{गुहा}} — \emph{guhā} — \textquotedblleft{}a cave\textquotedblright{}.

\begin{grammarlens}[title={What you've built}]
One more word for this chapter.
\end{grammarlens}

\subsection*{Your turn}
\begin{itemize}
\raggedright
  \item \emph{Say it:} \emph{guhā}
  \item \emph{Say it:} \emph{guhā}, once more
  \item \emph{Say it:} say \emph{guhā}
  \item \emph{Recall:} say the Sanskrit for the ocean, then the Sanskrit for a stream, then say \emph{guhā} again
\end{itemize}

\begin{tcolorbox}[breakable,colback=teal!4,colframe=teal!35!black,title={Before you move on}]
What does \sk{गुहा} mean? (\textquotedblleft{}A cave\textquotedblright{}.) Say it once more.
\end{tcolorbox}
\section[(dialogue)]{Words from chapters 191-193 — a first pass}

\label{lesson:SA-R196-first-pass-words-from-chapters-191-193}



\begin{quote}
Say the words for a stream and a cave.
\end{quote}

\subsection*{Your turn}
Say these aloud:

\begin{itemize}
\raggedright
  \item an earthen dish, a pitcher, a mortar, a pestle, a sieve
  \item a necklace, a crown, a lower garment, clothing, a garment
  \item a purchase, a sale, an object, an item, a tool
\end{itemize}

\begin{tcolorbox}[breakable,colback=teal!4,colframe=teal!35!black,title={Before you move on}]
An earthen dish? (\textbf{\sk{शरावः}}, \emph{śarāvaḥ}.) A pitcher? (\textbf{\sk{कुम्भः}}, \emph{kumbhaḥ}.) A mortar? (\textbf{\sk{उलूखलम्}}, \emph{ulūkhalam}.) A pestle? (\textbf{\sk{मुसलम्}}, \emph{musalam}.) A sieve? (\textbf{\sk{चालनी}}, \emph{cālanī}.) A necklace? (\textbf{\sk{हारः}}, \emph{hāraḥ}.) A crown? (\textbf{\sk{मुकुटम्}}, \emph{mukuṭam}.) A lower garment? (\textbf{\sk{अधोवस्त्रम्}}, \emph{adhovastram}.) Clothing? (\textbf{\sk{परिधानम्}}, \emph{paridhānam}.) A garment? (\textbf{\sk{वसनम्}}, \emph{vasanam}.) A purchase? (\textbf{\sk{क्रयः}}, \emph{krayaḥ}.) A sale? (\textbf{\sk{विक्रयः}}, \emph{vikrayaḥ}.) An object? (\textbf{\sk{वस्तु}}, \emph{vastu}.) An item? (\textbf{\sk{पदार्थः}}, \emph{padārthaḥ}.) A tool? (\textbf{\sk{उपकरणम्}}, \emph{upakaraṇam}.)
\end{tcolorbox}
\section[(dialogue)]{Words from chapters 194-196 — a second pass}

\label{lesson:SA-R196-second-pass-words-from-chapters-194-196}



\begin{quote}
Say the words for a prize and an invitation.
\end{quote}

\subsection*{Your turn}
Say these aloud:

\begin{itemize}
\raggedright
  \item a prize, an invitation, an assembly, a discussion group, a fair
  \item a port town, a capital, a district, a province, a boundary
  \item copper, a rock, the ocean, a stream, a cave
\end{itemize}

\begin{tcolorbox}[breakable,colback=teal!4,colframe=teal!35!black,title={Before you move on}]
A prize? (\textbf{\sk{पारितोषिकम्}}, \emph{pāritoṣikam}.) An invitation? (\textbf{\sk{निमन्त्रणम्}}, \emph{nimantraṇam}.) An assembly? (\textbf{\sk{सभा}}, \emph{sabhā}.) A discussion group? (\textbf{\sk{गोष्ठी}}, \emph{goṣṭhī}.) A fair? (\textbf{\sk{मेलकः}}, \emph{melakaḥ}.) A port town? (\textbf{\sk{पत्तनम्}}, \emph{pattanam}.) A capital? (\textbf{\sk{राजधानी}}, \emph{rājadhānī}.) A district? (\textbf{\sk{जनपदः}}, \emph{janapadaḥ}.) A province? (\textbf{\sk{प्रान्तः}}, \emph{prāntaḥ}.) A boundary? (\textbf{\sk{सीमा}}, \emph{sīmā}.) Copper? (\textbf{\sk{ताम्रम्}}, \emph{tāmram}.) A rock? (\textbf{\sk{पाषाणः}}, \emph{pāṣāṇaḥ}.) The ocean? (\textbf{\sk{रत्नाकरः}}, \emph{ratnākaraḥ}.) A stream? (\textbf{\sk{सरित्}}, \emph{sarit}.) A cave? (\textbf{\sk{गुहा}}, \emph{guhā}.)
\end{tcolorbox}
